%==============================================================================
\section{Supervised Learning --- Học máy có giám sát}
\label{sec:supervised-learning}
%==============================================================================

\subsection{Supervised Machine Learning là gì?}

\begin{dinhnghia}[title={Học có giám sát}]
Supervised learning huấn luyện mô hình bằng dataset \textbf{có nhãn}: mỗi mẫu gồm feature (đầu vào) và target (đầu ra) tương ứng.
target: học quan hệ giữa đầu vào và đầu ra sau nhiều lần huấn luyện.
\end{dinhnghia}

\subsection{Cách Supervised Learning hoạt động}

\begin{phuongphap}[title={Quy trình huấn luyện}]
\begin{itemize}
  \item Mô hình huấn luyện trên tập cặp đầu vào--đầu ra.
  \item Thuật toán phân tích dataset và học quan hệ giữa feature và nhãn/target đúng.
  \item Trong huấn luyện, mô hình ước lượng tham số bằng cách \textbf{tối thiểu hoá hàm mất mát} (loss): đo sai khác giữa dự đoán và giá trị thực.
  \item Cập nhật tham số lặp đến khi loss/error đủ nhỏ.
  \item Sau huấn luyện, tham số tối ưu; mô hình dự đoán được cho dữ liệu mới chưa thấy.
\end{itemize}
\end{phuongphap}

\subsection{Hai loại bài toán Supervised}

\subsubsection{Classification (phân loại)}

\begin{ynghia}[title={target}]
Dự đoán nhãn phân loại (đúng/sai, nam/nữ, có/không, \ldots) --- giá trị rời rạc, không có thứ tự bắt buộc.
\end{ynghia}

\begin{phuongphap}[title={Thuật toán phổ biến}]
decision trees, random forests, SVM, logistic regression, \ldots
\end{phuongphap}

\subsubsection{Regression (hồi quy)}

\begin{ynghia}[title={target}]
Dự đoán giá trị số \textbf{liên tục}; học quan hệ giữa biến độc lập và biến phụ thuộc.
\end{ynghia}

\begin{phuongphap}[title={Thuật toán phổ biến}]
linear regression, polynomial regression, Lasso regression, \ldots
\end{phuongphap}

\subsection{Thuật toán Supervised Learning}

\subsubsection{Linear Regression}

\begin{ynghia}[title={Hồi quy tuyến tính}]
Tìm quan hệ \textbf{tuyến tính} giữa feature và giá trị đầu ra để dự báo sự kiện tương lai.
Ứng dụng: phân tích cổ phiếu, dự báo thời tiết, \ldots
\end{ynghia}

\subsubsection{K-Nearest Neighbors (kNN)}

\begin{ynghia}[title={kNN}]
Kỹ thuật thống kê cho classification và regression: phân loại/dự đoán mẫu mới bằng khoảng cách tới các điểm trong tập huấn luyện.
\end{ynghia}

\begin{vidu}[title={Phân loại đối tượng lạ bằng kNN}]
\tsimg{11_mo_hinh_va_cac_kieu_hoc/nearest_neighbours.jpg}

Ba loại đối tượng (đỏ, xanh dương, xanh lá). Sau kNN, biên từng lớp như sau:

\tsimg{11_mo_hinh_va_cac_kieu_hoc/dataset_boundaries.jpg}

Điểm mới (chưa biết lớp):

\tsimg{11_mo_hinh_va_cac_kieu_hoc/depicted_figure.jpg}

Trực quan, điểm lạ thuộc lớp xanh dương.

Toán học: đo khoảng cách tới mọi điểm khác --- hầu hết láng giềng gần nhất là xanh dương; khoảng cách trung bình tới đỏ/xanh lá lớn hơn tới xanh dương $\Rightarrow$ phân lớp xanh dương.
kNN cũng dùng cho regression; có sẵn trong hầu hết thư viện ML.
\end{vidu}

\subsubsection{Decision Trees}

\begin{ynghia}[title={Cây quyết định}]
Cấu trúc dạng cây để ra quyết định và phân tích hậu quả.
Chia dữ liệu thành tập con theo feature; node trong là quyết định, lá là dự đoán cuối.
\end{ynghia}

\begin{vidu}[title={Ví dụ cây quyết định đơn giản (email)}]
\tsimg{11_mo_hinh_va_cac_kieu_hoc/flowchart_format.jpg}
Sơ đồ minh hoạ quy tắc đọc email (trivial trong ví dụ nhỏ).
Trong thực tế cây có thể lớn và phức tạp; cần nắm kỹ thuật tạo và duyệt cây.
\end{vidu}

\subsubsection{Naive Bayes}

\begin{ynghia}[title={Naive Bayes}]
Dùng xây classifier.
Ví dụ phân loại trái cây: dùng màu, kích thước, hình dạng; tính xác suất từng feature thỏa ràng buộc (ví dụ ``tròn, $\sim$10\,cm $\Rightarrow$ táo'').
Kết hợp xác suất các feature $\Rightarrow$ xác suất mẫu là táo.
Thường cần ít dữ liệu huấn luyện.
\end{ynghia}

\subsubsection{Logistic Regression}

\begin{ynghia}[title={Hồi quy logistic}]
Thuật toán thống kê ước lượng \textbf{xác suất} xảy ra sự kiện.
\end{ynghia}

\begin{vidu}[title={Phân tách lớp bằng đường biên}]
\tsimg{11_mo_hinh_va_cac_kieu_hoc/distribution_data_points.jpg}
Phân bố điểm đỏ và xanh; có thể vẽ đường biên --- điểm mới thuộc phía nào của đường quyết định thuộc lớp đó.
\end{vidu}

\subsubsection{Support Vector Machines (SVM)}

\begin{ynghia}[title={SVM}]
Dùng cho classification và regression.
Classification: tạo \textbf{siêu phẳng} tách lớp.
Regression: khớp đường hồi quy với sai số nhỏ.
\end{ynghia}

\begin{vidu}[title={Dữ liệu không tách tuyến tính}]
\tsimg{11_mo_hinh_va_cac_kieu_hoc/non_linear.jpg}
Ba lớp không tách được bằng đường thẳng; biên phi tuyến --- SVM hữu ích xác định ranh giới.
\end{vidu}

\subsubsection{Random Forest}

\begin{ynghia}[title={Random Forest}]
Thuật toán supervised linh hoạt cho classification và regression.
Kết hợp nhiều decision tree để cải thiện độ chính xác dự đoán.
\end{ynghia}

\begin{vidu}[title={Cách Random Forest hoạt động}]
\tsimg{11_mo_hinh_va_cac_kieu_hoc/random_forest_algorithm.jpg}
\end{vidu}

\subsubsection{Gradient Boosting}

\begin{ynghia}[title={Gradient Boosting}]
Kết hợp các weak learner (decision tree) thành mô hình mạnh; mô hình mới sửa lỗi mô hình trước.
target: tối thiểu hoá loss.
Dùng hiệu quả cho classification và regression.
\end{ynghia}

\subsection{Ưu điểm Supervised Learning}

\begin{tinhchat}[title={Ưu điểm}]
\begin{itemize}
  \item target rõ ràng $\Rightarrow$ cải thiện độ chính xác dự đoán.
  \item Dataset có nhãn giúp dự đoán và phân loại hiệu quả.
  \item Linh hoạt: spam detection, giá cổ phiếu, \ldots
\end{itemize}
\end{tinhchat}

\subsection{Nhược điểm}

\begin{luuy}[title={Nhược điểm}]
\begin{itemize}
  \item Cần \textbf{lượng lớn} dữ liệu có nhãn --- tốn kém, tốn thời gian thu thập và gán nhãn.
  \item Khó dự đoán chính xác nếu dữ liệu test khác phân phối train.
  \item Gán nhãn chính xác phức tạp, cần chuyên môn.
\end{itemize}
\end{luuy}

\subsection{Ứng dụng}

\begin{phuongphap}[title={Lĩnh vực ứng dụng}]
\textbf{Image recognition}: huấn luyện trên ảnh có nhãn; nhận diện khuôn mặt, phát hiện đối tượng.

\textbf{Predictive analytics}: học từ dữ liệu lịch sử có nhãn để nhận diện xu hướng, hỗ trợ quyết định kinh doanh dựa trên dữ liệu.
\end{phuongphap}
